\chapter{Web 1.0}\label{ch:ch1}

\section{Концептуальное назначение}\label{sec:ch1/sec1}

\[
--
\]

\section{Обобщенная структура}\label{sec:ch1/sec2}

%Сошлёмся на библиографию.
%Одна ссылка: \cite[с.~54]{Sokolov}\cite[с.~36]{Gaidaenko}.
%Две ссылки: \cite{Sokolov,Gaidaenko}.
%Ссылка на собственные работы: \cite{vakbib1, confbib2}.
%Много ссылок: %\cite[с.~54]{Lermontov,Management,Borozda} % такой «фокус»
%%вызывает biblatex warning относительно опции sortcites, потому что неясно, к
%%какому источнику относится уточнение о страницах, а bibtex об этой проблеме
%%даже не предупреждает
%\cite{Lermontov, Management, Borozda, Marketing, Constitution, FamilyCode,
%    Gost.7.0.53, Razumovski, Lagkueva, Pokrovski, Methodology, Berestova,
%    Kriger}%
%\ifnumequal{\value{bibliosel}}{0}{% Примеры для bibtex8
%    \cite{Sirotko, Lukina, Encyclopedia, Nasirova}%
%}{% Примеры для biblatex через движок biber
%    \cite{Sirotko2, Lukina2, Encyclopedia2, Nasirova2}%
%}%
%.
%И~ещё немного ссылок:~\cite{Article,Book,Booklet,Conference,Inbook,Incollection,Manual,Mastersthesis,
%    Misc,Phdthesis,Proceedings,Techreport,Unpublished}
%% Следует обратить внимание, что пробел после запятой внутри \cite{}
%% обрабатывается ожидаемо, а пробел перед запятой, может вызывать проблемы при
%% обработке ссылок.
%\cite{medvedev2006jelektronnye, CEAT:CEAT581, doi:10.1080/01932691.2010.513279,
%    Gosele1999161,Li2007StressAnalysis, Shoji199895, test:eisner-sample,
%    test:eisner-sample-shorted, AB_patent_Pomerantz_1968, iofis_patent1960}%
%\ifnumequal{\value{bibliosel}}{0}{% Примеры для bibtex8
%}{% Примеры для biblatex через движок biber
%    \cite{patent2h, patent3h, patent2}%
%}%
%.
%
%\ifnumequal{\value{bibliosel}}{0}{% Примеры для bibtex8
%Попытка реализовать несколько ссылок на конкретные страницы
%для \texttt{bibtex} реализации библиографии:
%[\citenum{Sokolov}, с.~54; \citenum{Gaidaenko}, с.~36].
%}{% Примеры для biblatex через движок biber
%Несколько источников (мультицитата):
%% Тут специально написано по-разному тире, для демонстрации, что
%% применение специальных тире в настоящий момент в biblatex приводит к непоказу
%% "с.".
%\cites[vii--x, 5, 7]{Sokolov}[v"--~x, 25, 526]{Gaidaenko}[vii--x, 5, 7]{Techreport},
%работает только в \texttt{biblatex} реализации библиографии.
%}%
%
%Ссылки на собственные работы:~\cite{vakbib1, confbib1}.
%
%%Сошлёмся на приложения: Приложение~\cref{app:A}, Приложение~\cref{app:B2}.
%
%Сошлёмся на формулу: формула~\cref{eq:equation1}.
%
%Сошлёмся на изображение: рисунок~\cref{fig:knuth}.
%
%Стандартной практикой является добавление к ссылкам префикса, характеризующего тип элемента.
%Это не является строгим требованием, но~позволяет лучше ориентироваться в документах большого размера.
%Например, для ссылок на~рисунки используется префикс \textit{fig},
%для ссылки на~таблицу "--- \textit{tab}.
%
%%В таблице \cref{tab:tab_pref} приложения~\cref{app:B4} приведён список рекомендуемых
%к использованию стандартных префиксов.
%
%В некоторых ситуациях возникает необходимость отойти от требований ГОСТ по оформлению ссылок на
%литературу.
%В таком случае можно воспользоваться дополнительными опциями пакета \verb+biblatex+.
%
%Например, в ссылке на книгу~\cite{sobenin_kdv} использование опции \verb+maxnames=4+ позволяет
%вывести имена всех четырёх авторов.
%По ГОСТ имена последних трёх авторов опускаются.
%
%Кроме того, часто возникают проблемы с транслитерованными инициалами. Некоторые буквы русского
%алфавита по правилам транслитерации записываются двумя буквами латинского алфавита (ю-yu, ё-yo и
%т.д.).
%Такие инициалы \verb+biblatex+ будет сокращать до одной буквы, что неверно.
%Поправить его работу можно использовав опцию \verb+giveninits=false+.
%Пример использования этой опции можно видеть в ссылке~\cite{initials}.

\subsection{Анализ топологии крупных сегментов Веба с помощью модели Бродера “Bow-tie”}\label{subsec:ch1/sec2/sub1}

\textbf{Введение.} В настоящее время все больше организаций, стремясь отразить свою деятельность в Вебе, создают отдельные веб-ресурсы (сайты, группы сайтов) и заботятся об улучшении их рейтингов в информационно-поисковых системах с целью накопления символического капитала, повышения прямых продаж. Важную роль в формировании рейтинга сайтов и локализированных веб-сегментов в целом играют поисковые машины.

С 2006 года основные поисковые машины (например, в мире -- Google, в России -- Яндекс) в своих методах ранжирования наряду с классическими факторами \cite{Kleinberg, BrinPage, Chakrabarti} (частота встречаемости слов, объем цитирования и авторитетность сайтов, частота обновления сайтов и др.) по-новому используют факторы, связанные с поведением пользователя на веб-ресурсах \cite{GuhaKunduBhadra,AntoniouPlegasTsakalidis,FeuerSavevAslam}. Так, действия пользователя, ранее приводившие к повышению рейтингов сайтов (например, частота захода пользователей), ныне могут приводить к “пессимизации” показателей, поскольку поисковые системы становятся более “социальными” и учитывают уже не только профиль пользователя и частоту захода на ресурсы, но также мотивацию посетителей и стратегии их поведения (частоту повторного захода, время нахождения на странице и сайте, логику и маршруты переходов, пользовательские интересы, тип потребляемого контента и многие иные факторы). При ранжировании современные поисковые машины в совокупности учитывают около 300-500 таких факторов.

С другой стороны, существует небольшое количество структурных характеристик Веб-сегментов \cite{ChoRoy}, которые существенно влияют практически на все эти факторы. Имеется значительное число исследований \cite{Kleinberg,ChoRoy,BroderKumarMaghoul,AguilloGranadinoOrtega,StuartThelwallHarries,Chakrabarti,Thelwall,Pechnikov,PechnikovNwohiri}, в которых строится модель фрагментов веба, описывающаяся сравнительно небольшим набором параметров.

Общим направлением исследований авторов является изучение связей между рейтингом сайта и его глобальными характеристиками. По-видимому, эти связи различны для сайтов разных категорий, например, естественнонаучной направленности и гуманитарных. Для выявления существенных различий между категориями сайтов, необходимо исследовать большое количество сайтов с последующей статистической обработкой материала. Авторы располагают разработанной ими аналитической системой для вебометрических исследований \cite{BlekanovSergeevMaksimov,BlekanovSergeevMartynenko}, способной выполнить эту работу. Ввиду чрезвычайно больших затрат машинных ресурсов на подобное исследование, представляет интерес машинный эксперимент, заключающийся в определении затрат на полное исследование типичных сайтов, описание которого и является темой данной статьи.

\textbf{Основная часть.} В эксперименте ставилась задача построения топологии на основе модели Бродера “Bow-tie” \cite{BroderKumarMaghoul,Thelwall} для типичных крупных гуманитарно-ориентированных и естественнонаучно-ориентированных сегментов Веба с последующим вычислением их структурных характеристик и определением произведенных затрат ресурсов.

В качестве гуманитарно-ориентированного веб-сегмента был выбран сайт Факультета журналистики (JF) Санкт-Петербургского государственного университета (СПбГУ) -- www.jf.spbu.ru, естественнонаучно-ориентированного -- сайт факультета Прикладной математики -- процессов управления (APMATH) СПбГУ -- www.apmath.spbu.ru.

С помощью специализированной сборки аналитической системы для вебометрических исследований на основе ядра поискового робота, успешно апробированной в исследованиях \cite{BlekanovSergeev,MaksimovBlekanov} , были выявлены гиперссылочные структуры обоих сайтов, которые описываются в виде двух ориентированных веб-графов \(G_{apmath}\) и \(G_{jf}\), в которых вершины соответствуют страницам, а дуги – соединяющим эти страницы гиперссылкам. Для получения глобальных характеристик использовалась модель Бродера “Bow-tie” \cite{BroderKumarMaghoul} , которая во множестве вершин исследуемого графа выделяет четыре основных подмножества: 1) центральное ядро (компонента сильной связности SCC); 2) истоки (IN); 3) стоки (OUT); 4) “TENDRILS” и “TUBES”. Поиск сильно-связной компоненты выполнялся с помощью алгоритма Косарайю \cite{Sedgewick} (Kosaraju's algorithm), время выполнения которого для разреженных графов (а такими являются исследуемые веб-графы) пропорционально сумме всех их вершин и направленных ребер.

В эксперименте вычислялись следующие характеристики веб-сайтов:
\begin{itemize}
	\item размер центрального ядра \(S_{scc}\);
	\item размер стоков \(S_{out}\) и истоков \(S_{in}\);
	\item размер “TENDRILS” и “TUBES” \(S_{tubes}\);
	\item отношение \(|{S_{scc}| + |S_{in}|}\) к \(|S_{scc}| + |S_{out}|\).
\end{itemize}

Кроме того, для каждой страницы центрального ядра был введен показатель ее связности с другими страницами из ядра:
\[
	\text{\textit{VC}}(u, v) = 
	\begin{cases} 
		1, & \text{\textit{if }} u \rightarrow v,\\
		0, & \text{\textit{else}}
	\end{cases},
	\forall u, v \in S_{scc}
\]

Этот показатель лег в основу усредненной меры связности каждой страницы ядра \(\text{\textit{MA\text{\textit{VC}}}}(u)\) и усредненной меры связности центрального ядра \(\text{\textit{MAC}}(S_{scc})\):
\[
	\text{\textit{MAVC}}(u) = \frac{1}{\lvert S_{scc} \rvert - 1} \sum_{\forall v \in S_{scc}, \; v \neq u}\text{\textit{VC}}(u, v),
\] 
\[
	\text{\textit{MAC}}(S_{scc}) = \frac{1}{|S_{scc}|} \sum_{\forall u \in S_{scc}}\text{\textit{MAVC}}(u).
\]

Приведенные меры были введены для оценки качества связанности каждой топологии, полученной в ходе поставленного эксперимента.

Для оценки трудоемкости эксперимента по каждому исследуемому сайту использовались следующие характеристики:
\begin{itemize}
	\item время (\(T_{total}^{js}\) и \(T_{total}^{apmath}\)), необходимое для выявления топологии и вычисления указанных обобщенных характеристик;
	\item объем памяти (\(\text{\textit{VMem}}_{js}\) и \(\text{\textit{VMem}}_{apmath}\)), требуемый для хранения обработанных данных;
	\item объем Веб трафика \(\text{\textit{Traffic}}_{js}\)  и \(\text{\textit{Traffic}}_{apmath}\)).
\end{itemize}

\textbf{Результаты эксперимента.} Эксперимент проводился на тестовой персональной электронно-вычислительной машине (ПЭВМ) с процессором \textit{Intel Core i7 CPU 950 @ 3.07 GHz x 8}, \textit{12 GB} объемом жесткого диска и операционной системой \textit{Ubuntu 13.10}. Загрузка данных с Интернета выполнялась посредством канала с пропускной способностью до \textit{50 Мбит/с}.

В ходе эксперимента было выявлено, что Веб-граф \(G_{apmath}\) содержит \textit{26 148} вершин и \textit{2 025 909} направленных ребер, а Веб-граф \(G_{jf}\) -- \textit{27 361} вершин и \textit{4 252 424} направленных ребра. На  Рис.~\cref{fig:webGraphs} приведены топологии, построенные для каждого полученного графа с помощью модели Бродера “Bow-tie”.

\begin{figure}[ht]
    \centerfloat{
	        \hfill
	        \subcaptionbox[List-of-Figures entry]{\label{fig:webGraphs-1}}{%
		            \includegraphics[width=0.5\linewidth]{webGraph(apmath)}}
%	        \hfill
	        \subcaptionbox{\label{fig:webGraphs-2}}{%
		            \includegraphics[width=0.5\linewidth]{webGraph(jf)}}
	        \hfill
	    }
    \caption{Топологии Веб-графов \(G_{apmath}\) и \(G_{jf}\), соответственно.}\label{fig:webGraphs}
\end{figure}

Значительную долю в топологии Веб-графа \(G_{apmath}\) среди всех ее компонент занимает SCC (\textit{65.56\%} страниц от общего их числа), \textit{8.8\%} занимают истоки (IN), - \textit{25.5\%} стоки (OUT), \textit{3.96\%} -- “TENDRILS” и “TUBES” (Рис.~\cref{fig:webGraphs-1}). В то время как доли компонент топологии Веб-графа (Рис.~\cref{fig:webGraphs}\subcaptionref*{fig:webGraphs-2}) приблизительно одинаковые (более \textit{46\%} от общего числа всех страниц) за исключением компоненты OUT (\textit{0.76\%}).

В Таблице \cref{tab:webGraphTable} приведены полученные численные значения структурных характеристик, которые были введены в эксперименте для анализа исследуемых сайтов.

\begin{table} [htbp]%
    \centering
    \caption{Структурные характеристики Веб-графов \(G_{apmath}\) и \(G_{jf}\).}%
    \label{tab:webGraphTable}% label всегда желательно идти после caption
    \renewcommand{\arraystretch}{1.5}%% Увеличение расстояния между рядами, для улучшения восприятия.
    \begin{SingleSpace}
    	\begin{tabulary}{\textwidth}{@{}>{\zz}L >{\zz}C >{\zz}C@{}} %Вертикальные полосы не используются принципиально, как и лишние горизонтальные (допускается по ГОСТ 2.105 пункт 4.4.5) % @{} позволяет прижиматься к краям
		            \toprule     %%% верхняя линейка
		            Характеристика & Топология сайта APMATH &Топология сайта JS \\
		            \midrule %%% тонкий разделитель. Отделяет названия столбцов. Обязателен по ГОСТ 2.105 пункт 4.4.5
		            Мощность множества ядра \(\lvert S_{scc} \rvert\)  & 17142 & 13893    \\
		            Мощность множества IN \(\lvert S_{in} \rvert\)       & 2302     & 13165    \\
		            Мощность множества OUT \(\lvert S_{out} \rvert\)       & 6669     & 206    \\
					Мощность множества \newline TENDRILS и TUBES  \(\lvert S_{tubes} \rvert\)      & 1035    & 12809     \\
					Отношение \(\frac{\lvert S_{scc} \rvert + \lvert S_{in} \rvert}{\lvert S_{scc} \rvert + \lvert S_{out} \rvert}\) & 0.82 & 1.92 \\
					Усредненная мера связности центрального ядра \(\text{\textit{MAC}}(S_{scc})\) & 0.00244 & 0.00193 \\
					Усредненная мера связности всего Веб-графа \textit{MAC(G)} & 0.00106 & 0.00098 \\
		            \bottomrule %%% нижняя линейка
		        \end{tabulary}%
	    \end{SingleSpace}
\end{table}

Затраты ресурсов, требуемых для полного анализа исследуемых в эксперименте сегментов Веба, приведены в Таблице \cref{tab:webGraphCostTable}.

\begin{table}[ht]%
	\caption{Характеристики, оценивающие трудоемкость поставленного эксперимента.}%
	\label{tab:webGraphCostTable}% label всегда желательно идти после caption
    \renewcommand{\arraystretch}{1.6}%% Увеличение расстояния между рядами, для улучшения восприятия.
    \def\tabularxcolumn#1{m{#1}}
    \begin{tabularx}{\textwidth}{@{}>{\raggedright}X>{\centering}m{3.5cm}  >{\centering\arraybackslash}m{3.5cm}@{}}% Вертикальные полосы не используются принципиально, как и лишние горизонтальные (допускается по ГОСТ 2.105 пункт 4.4.5) % @{} позволяет прижиматься к краям
			\toprule     %%% верхняя линейка
			Характеристика & Сайт APMATH & Сайт JS \\
			\midrule %%% тонкий разделитель. Отделяет названия столбцов. Обязателен по ГОСТ 2.105 пункт 4.4.5
			Время, необходимое для построения топологии и вычисления структурных характеристик, \textit{min.}  & \(T_{total}^{apmath} = 672\) &  \(T_{total}^{js} = 271\)  \\
			Среднее время загрузки одной страницы, \textit{min.} & 0.026 & 0.001 \\
			Объем памяти на жестком диске для хранения данных, \textit{MB} & \(\text{\textit{VMem}}_{apmath} = 154.743\) & \(\text{\textit{VMem}}_{js} = 318.463\) \\
			Минимальный объем \newline оперативной памяти \newline для обработки данных, \textit{GB} & \(\approx 1.5\) & \(\approx 1.5\) \\
			Средний объем памяти, требуемый  для полной обработки \newline одной веб-страницы, \textit{kB} & 6.06 & 11.92 \\
			Объем Веб трафика, \textit{GB} &  \(\text{\textit{Traffic}}_{apmath} = 9.66\) &  \(\text{\textit{Traffic}}_{js} = 3.1\) \\
			\bottomrule %%% нижняя линейка
	    \end{tabularx}%
\end{table}

\textbf{Выводы}. Результаты эксперимента подтвердили существенные различия между гиперссылочной структурой представителя сайтов естественнонаучной направленности и структурой представителя гуманитарно-ориентированных сайтов. Так, например, в топологии первого сайта (Рис.~\cref{fig:webGraphs-1}) большая доля всех веб-страниц сосредоточена в центральном ядре, которое имеет большую связность между своими элементами по сравнению с элементами ядра второго сайта (Рис.~\cref{fig:webGraphs}\subcaptionref*{fig:webGraphs-2}). Кроме, того значительное различие компоненты OUT первой топологии в сравнении со второй объясняется наличием в ней большого числа полнотекстовых документов (в формате PDF, DOC, DOCX и др.). Однако, количество гиперссылок в Веб-графе \(G_{jf}\) и характеристик размеров компонент SCC, IN, TUBES/TENDRILS (Таб.~\cref{tab:webGraphTable}) говорят о хорошей коммуникабельности между компонентами топологии Веб-графа \(G_{jf}\).

Для получения статистически достоверных результатов необходимо провести экспериментальные исследования топологий большого числа сайтов для каждой выделенной категории. Опираясь на полученные значения характеристик (Таб.~\cref{tab:webGraphCostTable}), оценивающие трудоемкость поставленного в данной статье эксперимента, можно приблизительно оценить ресурсные затраты на полное исследование любого количества сайтов.

\section{Сбор данных}\label{sec:ch1/sec3}

%Благодаря пакету \textit{icomma}, \LaTeX~одинаково хорошо воспринимает
%в~качестве десятичного разделителя и запятую (\(3,1415\)), и точку (\(3.1415\)).

\subsection{Построение тематико-ориентированных веб-краулеров с использованием обобщенного ядра}\label{subsec:ch1/sec3/sub1}

%Вот так может выглядеть формула, которую необходимо вставить в~строку
%по~тексту: \(x \approx \sin x\) при \(x \to 0\).
%
%А вот так выглядит ненумерованная отдельностоящая формула c подстрочными
%и надстрочными индексами:
%\[
%    (x_1+x_2)^2 = x_1^2 + 2 x_1 x_2 + x_2^2
%\]
%
%Формула с неопределенным интегралом:
%\[
%    \int f(\alpha+x)=\sum\beta
%\]
%
%При использовании дробей формулы могут получаться очень высокие:
%\[
%    \frac{1}{\sqrt{2}+
%        \displaystyle\frac{1}{\sqrt{2}+
%            \displaystyle\frac{1}{\sqrt{2}+\cdots}}}
%\]
%
%В формулах можно использовать греческие буквы:
%%Все \original... команды заранее, ради этого примера, определены в Dissertation\userstyles.tex
%\[
%    \alpha\beta\gamma\delta\originalepsilon\epsilon\zeta\eta\theta%
%    \vartheta\iota\kappa\varkappa\lambda\mu\nu\xi\pi\varpi\rho\varrho%
%    \sigma\varsigma\tau\upsilon\originalphi\phi\chi\psi\omega\Gamma\Delta%
%    \Theta\Lambda\Xi\Pi\Sigma\Upsilon\Phi\Psi\Omega
%\]
%\[%https://texfaq.org/FAQ-boldgreek
%    \boldsymbol{\alpha\beta\gamma\delta\originalepsilon\epsilon\zeta\eta%
%        \theta\vartheta\iota\kappa\varkappa\lambda\mu\nu\xi\pi\varpi\rho%
%        \varrho\sigma\varsigma\tau\upsilon\originalphi\phi\chi\psi\omega\Gamma%
%        \Delta\Theta\Lambda\Xi\Pi\Sigma\Upsilon\Phi\Psi\Omega}
%\]
%
%Для добавления формул можно использовать пары \verb+$+\dots\verb+$+ и \verb+$$+\dots\verb+$$+,
%но~они считаются устаревшими.
%Лучше использовать их функциональные аналоги \verb+\(+\dots\verb+\)+ и \verb+\[+\dots\verb+\]+.

\subsection{Веб-краулер как инструмент для вебометрических исследований на примере анализа Веб-пространства СПбГУ}\label{subsec:ch1/sec3/sub2}

%Вот так можно написать две формулы, не нумеруя их, чтобы знаки <<равно>> были
%строго друг под другом:
%\begin{align}
%    f_W & =  \min \left( 1, \max \left( 0, \frac{W_{soil} / W_{max}}{W_{crit}} \right)  \right), \nonumber \\
%    f_T & =  \min \left( 1, \max \left( 0, \frac{T_s / T_{melt}}{T_{crit}} \right)  \right), \nonumber
%\end{align}
%
%Выровнять систему ещё и по переменной \( x \) можно, используя окружение
%\verb|alignedat| из пакета \verb|amsmath|. Вот так:
%\[
%|x| = \left\{
%\begin{alignedat}{2}
%    &&x, \quad &\text{eсли } x\geqslant 0 \\
%    &-&x, \quad & \text{eсли } x<0
%\end{alignedat}
%\right.
%\]
%Здесь первый амперсанд (в исходном \LaTeX\ описании формулы) означает
%выравнивание по~левому краю, второй "--- по~\( x \), а~третий "--- по~слову
%<<если>>. Команда \verb|\quad| делает большой горизонтальный пробел.
%
%Ещё вариант:
%\[
%    |x|=
%    \begin{cases}
%        \phantom{-}x, \text{если } x \geqslant 0 \\
%        -x, \text{если } x<0
%    \end{cases}
%\]
%
%Кроме того, для  нумерованных формул \verb|alignedat| делает вертикальное
%выравнивание номера формулы по центру формулы. Например, выравнивание
%компонент вектора:
%\begin{equation}
%    \label{eq:2p3}
%    \begin{alignedat}{2}
%        {\mathbf{N}}_{o1n}^{(j)} = \,{\sin} \phi\,n\!\left(n+1\right)
%        {\sin}\theta\,
%        \pi_n\!\left({\cos} \theta\right)
%        \frac{
%        z_n^{(j)}\!\left( \rho \right)
%        }{\rho}\,
%        &{\boldsymbol{\hat{\mathrm e}}}_{r}\,+   \\
%        +\,
%        {\sin} \phi\,
%        \tau_n\!\left({\cos} \theta\right)
%        \frac{
%        \left[\rho z_n^{(j)}\!\left( \rho \right)\right]^{\prime}
%        }{\rho}\,
%        &{\boldsymbol{\hat{\mathrm e}}}_{\theta}\,+   \\
%        +\,
%        {\cos} \phi\,
%        \pi_n\!\left({\cos} \theta\right)
%        \frac{
%        \left[\rho z_n^{(j)}\!\left( \rho \right)\right]^{\prime}
%        }{\rho}\,
%        &{\boldsymbol{\hat{\mathrm e}}}_{\phi}\:.
%    \end{alignedat}
%\end{equation}
%
%Ещё об отступах. Иногда для лучшей <<читаемости>> формул полезно
%немного исправить стандартные интервалы \LaTeX\ с учётом логической
%структуры самой формулы. Например в формуле~\cref{eq:2p3} добавлен
%небольшой отступ \verb+\,+ между основными сомножителями, ниже
%результат применения всех вариантов отступа:
%\begin{align*}
%    \backslash!             & \quad f(x) = x^2\! +3x\! +2         \\
%    \mbox{по-умолчанию}     & \quad f(x) = x^2+3x+2               \\
%    \backslash,             & \quad f(x) = x^2\, +3x\, +2         \\
%    \backslash{:}           & \quad f(x) = x^2\: +3x\: +2         \\
%    \backslash;             & \quad f(x) = x^2\; +3x\; +2         \\
%    \backslash \mbox{space} & \quad f(x) = x^2\ +3x\ +2           \\
%    \backslash \mbox{quad}  & \quad f(x) = x^2\quad +3x\quad +2   \\
%    \backslash \mbox{qquad} & \quad f(x) = x^2\qquad +3x\qquad +2
%\end{align*}
%
%Можно использовать разные математические алфавиты:
%\begin{align}
%    \mathcal{ABCDEFGHIJKLMNOPQRSTUVWXYZ} \nonumber  \\
%    \mathfrak{ABCDEFGHIJKLMNOPQRSTUVWXYZ} \nonumber \\
%    \mathbb{ABCDEFGHIJKLMNOPQRSTUVWXYZ} \nonumber
%\end{align}
%
%Посмотрим на систему уравнений на примере аттрактора Лоренца:
%
%\[
%\left\{
%\begin{array}{rl}
%    \dot x = & \sigma (y-x)  \\
%    \dot y = & x (r - z) - y \\
%    \dot z = & xy - bz
%\end{array}
%\right.
%\]
%
%А для вёрстки матриц удобно использовать многоточия:
%\[
%    \left(
%        \begin{array}{ccc}
%            a_{11} & \ldots & a_{1n} \\
%            \vdots & \ddots & \vdots \\
%            a_{n1} & \ldots & a_{nn} \\
%        \end{array}
%    \right)
%\]
%
%%\subsection{Нумерованные формулы}\label{subsec:ch1/sec3/sub3}
%
%А вот так пишется нумерованная формула:
%\begin{equation}
%    \label{eq:equation1}
%    e = \lim_{n \to \infty} \left( 1+\frac{1}{n} \right) ^n
%\end{equation}
%
%Нумерованных формул может быть несколько:
%\begin{equation}
%    \label{eq:equation2}
%    \lim_{n \to \infty} \sum_{k=1}^n \frac{1}{k^2} = \frac{\pi^2}{6}
%\end{equation}
%
%Впоследствии на формулы~\cref{eq:equation1, eq:equation2} можно ссылаться.
%
%Сделать так, чтобы номер формулы стоял напротив средней строки, можно,
%используя окружение \verb|multlined| (пакет \verb|mathtools|) вместо
%\verb|multline| внутри окружения \verb|equation|. Вот так:
%\begin{equation} % \tag{S} % tag - вписывает свой текст
%    \label{eq:equation3}
%    \begin{multlined}
%        1+ 2+3+4+5+6+7+\dots + \\
%        + 50+51+52+53+54+55+56+57 + \dots + \\
%        + 96+97+98+99+100=5050
%    \end{multlined}
%\end{equation}
%
%Уравнения~\cref{eq:subeq_1,eq:subeq_2} демонстрируют возможности
%окружения \verb|\subequations|.
%\begin{subequations}
%    \label{eq:subeq_1}
%    \begin{gather}
%        y = x^2 + 1 \label{eq:subeq_1-1} \\
%        y = 2 x^2 - x + 1 \label{eq:subeq_1-2}
%    \end{gather}
%\end{subequations}
%Ссылки на отдельные уравнения~\cref{eq:subeq_1-1,eq:subeq_1-2,eq:subeq_2-1}.
%\begin{subequations}
%    \label{eq:subeq_2}
%    \begin{align}
%        y & = x^3 + x^2 + x + 1 \label{eq:subeq_2-1} \\
%        y & = x^2
%    \end{align}
%\end{subequations}
%
%%\subsection{Форматирование чисел и размерностей величин}\label{sec:units}
%
%Числа форматируются при помощи команды \verb|\num|:
%\num{5,3};
%\num{2,3e8};
%\num{12345,67890};
%\num{2,6 d4};
%\num{1+-2i};
%\num{.3e45};
%\num[exponent-base=2]{5 e64};
%\num[exponent-base=2,exponent-to-prefix]{5 e64};
%\num{1.654 x 2.34 x 3.430}
%\num{1 2 x 3 / 4}.
%Для написания последовательности чисел можно использовать команды \verb|\numlist| и \verb|\numrange|:
%\numlist{10;30;50;70}; \numrange{10}{30}.
%Значения углов можно форматировать при помощи команды \verb|\ang|:
%\ang{2.67};
%\ang{30,3};
%\ang{-1;;};
%\ang{;-2;};
%\ang{;;-3};
%\ang{300;10;1}.
%
%Обратите внимание, что ГОСТ запрещает использование знака <<->> для обозначения отрицательных чисел
%за исключением формул, таблиц и~рисунков.
%Вместо него следует использовать слово <<минус>>.
%
%Размерности можно записывать при помощи команд \verb|\si| и \verb|\SI|:
%\si{\farad\squared\lumen\candela};
%\si{\joule\per\mole\per\kelvin};
%\si[per-mode = symbol-or-fraction]{\joule\per\mole\per\kelvin};
%\si{\metre\per\second\squared};
%\SI{0.10(5)}{\neper};
%\SI{1.2-3i e5}{\joule\per\mole\per\kelvin};
%\SIlist{1;2;3;4}{\tesla};
%\SIrange{50}{100}{\volt}.
%Список единиц измерений приведён в таблицах~\cref{tab:unit:base,
%    tab:unit:derived,tab:unit:accepted,tab:unit:physical,tab:unit:other}.
%Приставки единиц приведены в~таблице~\cref{tab:unit:prefix}.
%
%С дополнительными опциями форматирования можно ознакомиться в~описании пакета \texttt{siunitx};
%изменить или добавить единицы измерений можно в~файле \texttt{siunitx.cfg}.
%
%\begin{table}
%    \centering
%    \captionsetup{justification=centering} % выравнивание подписи по-центру
%    \caption{Основные величины СИ}\label{tab:unit:base}
%    \begin{tabular}{llc}
%        \toprule
%        Название  & Команда                 & Символ         \\
%        \midrule
%        Ампер     & \verb|\ampere| & \si{\ampere}   \\
%        Кандела   & \verb|\candela| & \si{\candela}  \\
%        Кельвин   & \verb|\kelvin| & \si{\kelvin}   \\
%        Килограмм & \verb|\kilogram| & \si{\kilogram} \\
%        Метр      & \verb|\metre| & \si{\metre}    \\
%        Моль      & \verb|\mole| & \si{\mole}     \\
%        Секунда   & \verb|\second| & \si{\second}   \\
%        \bottomrule
%    \end{tabular}
%\end{table}
%
%\begin{table}
%    \small
%    \centering
%    \begin{threeparttable}% выравнивание подписи по границам таблицы
%        \caption{Производные единицы СИ}\label{tab:unit:derived}
%        \begin{tabular}{llc|llc}
%            \toprule
%            Название       & Команда                 & Символ              & Название & Команда & Символ \\
%            \midrule
%            Беккерель      & \verb|\becquerel| & \si{\becquerel}     &
%            Ньютон         & \verb|\newton| & \si{\newton}                                      \\
%            Градус Цельсия & \verb|\degreeCelsius| & \si{\degreeCelsius} &
%            Ом             & \verb|\ohm| & \si{\ohm}                                         \\
%            Кулон          & \verb|\coulomb| & \si{\coulomb}       &
%            Паскаль        & \verb|\pascal| & \si{\pascal}                                      \\
%            Фарад          & \verb|\farad| & \si{\farad}         &
%            Радиан         & \verb|\radian| & \si{\radian}                                      \\
%            Грей           & \verb|\gray| & \si{\gray}          &
%            Сименс         & \verb|\siemens| & \si{\siemens}                                     \\
%            Герц           & \verb|\hertz| & \si{\hertz}         &
%            Зиверт         & \verb|\sievert| & \si{\sievert}                                     \\
%            Генри          & \verb|\henry| & \si{\henry}         &
%            Стерадиан      & \verb|\steradian| & \si{\steradian}                                   \\
%            Джоуль         & \verb|\joule| & \si{\joule}         &
%            Тесла          & \verb|\tesla| & \si{\tesla}                                       \\
%            Катал          & \verb|\katal| & \si{\katal}         &
%            Вольт          & \verb|\volt| & \si{\volt}                                        \\
%            Люмен          & \verb|\lumen| & \si{\lumen}         &
%            Ватт           & \verb|\watt| & \si{\watt}                                        \\
%            Люкс           & \verb|\lux| & \si{\lux}           &
%            Вебер          & \verb|\weber| & \si{\weber}                                       \\
%            \bottomrule
%        \end{tabular}
%    \end{threeparttable}
%\end{table}
%
%\begin{table}
%    \centering
%    \begin{threeparttable}% выравнивание подписи по границам таблицы
%        \caption{Внесистемные единицы}\label{tab:unit:accepted}
%
%        \begin{tabular}{llc}
%            \toprule
%            Название        & Команда                 & Символ          \\
%            \midrule
%            День            & \verb|\day| & \si{\day}       \\
%            Градус          & \verb|\degree| & \si{\degree}    \\
%            Гектар          & \verb|\hectare| & \si{\hectare}   \\
%            Час             & \verb|\hour| & \si{\hour}      \\
%            Литр            & \verb|\litre| & \si{\litre}     \\
%            Угловая минута  & \verb|\arcminute| & \si{\arcminute} \\
%            Угловая секунда & \verb|\arcsecond| & \si{\arcsecond} \\ %
%            Минута          & \verb|\minute| & \si{\minute}    \\
%            Тонна           & \verb|\tonne| & \si{\tonne}     \\
%            \bottomrule
%        \end{tabular}
%    \end{threeparttable}
%\end{table}
%
%\begin{table}
%    \centering
%    \captionsetup{justification=centering}
%    \caption{Внесистемные единицы, получаемые из эксперимента}\label{tab:unit:physical}
%    \begin{tabular}{llc}
%        \toprule
%        Название                & Команда                 & Символ                 \\
%        \midrule
%        Астрономическая единица & \verb|\astronomicalunit| & \si{\astronomicalunit} \\
%        Атомная единица массы   & \verb|\atomicmassunit| & \si{\atomicmassunit}   \\
%        Боровский радиус        & \verb|\bohr| & \si{\bohr}             \\
%        Скорость света          & \verb|\clight| & \si{\clight}           \\
%        Дальтон                 & \verb|\dalton| & \si{\dalton}           \\
%        Масса электрона         & \verb|\electronmass| & \si{\electronmass}     \\
%        Электрон Вольт          & \verb|\electronvolt| & \si{\electronvolt}     \\
%        Элементарный заряд      & \verb|\elementarycharge| & \si{\elementarycharge} \\
%        Энергия Хартри          & \verb|\hartree| & \si{\hartree}          \\
%        Постоянная Планка       & \verb|\planckbar| & \si{\planckbar}        \\
%        \bottomrule
%    \end{tabular}
%\end{table}
%
%\begin{table}
%    \centering
%    \begin{threeparttable}% выравнивание подписи по границам таблицы
%        \caption{Другие внесистемные единицы}\label{tab:unit:other}
%        \begin{tabular}{llc}
%            \toprule
%            Название                  & Команда                 & Символ             \\
%            \midrule
%            Ангстрем                  & \verb|\angstrom| & \si{\angstrom}     \\
%            Бар                       & \verb|\bar| & \si{\bar}          \\
%            Барн                      & \verb|\barn| & \si{\barn}         \\
%            Бел                       & \verb|\bel| & \si{\bel}          \\
%            Децибел                   & \verb|\decibel| & \si{\decibel}      \\
%            Узел                      & \verb|\knot| & \si{\knot}         \\
%            Миллиметр ртутного столба & \verb|\mmHg| & \si{\mmHg}         \\
%            Морская миля              & \verb|\nauticalmile| & \si{\nauticalmile} \\
%            Непер                     & \verb|\neper| & \si{\neper}        \\
%            \bottomrule
%        \end{tabular}
%    \end{threeparttable}
%\end{table}
%
%\begin{table}
%    \small
%    \centering
%    \begin{threeparttable}% выравнивание подписи по границам таблицы
%        \caption{Приставки СИ}\label{tab:unit:prefix}
%        \begin{tabular}{llcc|llcc}
%            \toprule
%            Приставка & Команда                  & Символ      & Степень &
%            Приставка & Команда                  & Символ      & Степень   \\
%            \midrule
%            Иокто     & \verb|\yocto|  & \si{\yocto} & -24     &
%            Дека      & \verb|\deca|  & \si{\deca}  & 1         \\
%            Зепто     & \verb|\zepto|  & \si{\zepto} & -21     &
%            Гекто     & \verb|\hecto|  & \si{\hecto} & 2         \\
%            Атто      & \verb|\atto|  & \si{\atto}  & -18     &
%            Кило      & \verb|\kilo|  & \si{\kilo}  & 3         \\
%            Фемто     & \verb|\femto|  & \si{\femto} & -15     &
%            Мега      & \verb|\mega|  & \si{\mega}  & 6         \\
%            Пико      & \verb|\pico|  & \si{\pico}  & -12     &
%            Гига      & \verb|\giga|  & \si{\giga}  & 9         \\
%            Нано      & \verb|\nano|  & \si{\nano}  & -9      &
%            Терра     & \verb|\tera|  & \si{\tera}  & 12        \\
%            Микро     & \verb|\micro|  & \si{\micro} & -6      &
%            Пета      & \verb|\peta|  & \si{\peta}  & 15        \\
%            Милли     & \verb|\milli|  & \si{\milli} & -3      &
%            Екса      & \verb|\exa|  & \si{\exa}   & 18        \\
%            Санти     & \verb|\centi|  & \si{\centi} & -2      &
%            Зетта     & \verb|\zetta|  & \si{\zetta} & 21        \\
%            Деци      & \verb|\deci| & \si{\deci}  & -1      &
%            Иотта     & \verb|\yotta| & \si{\yotta} & 24        \\
%            \bottomrule
%        \end{tabular}
%    \end{threeparttable}
%\end{table}

%\subsection{Заголовки с формулами: \texorpdfstring{\(a^2 + b^2 = c^2\)}{%
%        a\texttwosuperior\ + b\texttwosuperior\ = c\texttwosuperior},
%   \texorpdfstring{\(\left\vert\textrm{{Im}}\Sigma\left(
%            \protect\varepsilon\right)\right\vert\approx const\)}{|ImΣ (ε)| ≈ const},
%    \texorpdfstring{\(\sigma_{xx}^{(1)}\)}{σ\_\{xx\}\textasciicircum\{(1)\}}
%}\label{subsec:with_math}

%Пакет \texttt{hyperref} берёт текст для закладок в pdf-файле из~аргументов
%команд типа \verb|\section|, которые могут содержать математические формулы,
%а~также изменения цвета текста или шрифта, которые не отображаются в~закладках.
%Чтобы использование формул в заголовках не вызывало в~логе компиляции появление
%предупреждений типа <<\texttt{Token not allowed in~a~PDF string
%    (Unicode):(hyperref) removing...}>>, следует использовать конструкцию
%\verb|\texorpdfstring{}{}|, где в~первых фигурных скобках указывается
%формула, а~во~вторых "--- запись формулы для закладок.

\textbf{Введение.} За последние несколько лет в области информационного Веб-поиска все чаще появляются задачи, связанные с развивающимся научным направлением вебометрика (webometrics) \cite{HolmbergThelwall,Pechnikov,PechnikovChirkovChuiko,BlekanovSergeevPechnikov}. К актуальным направлениям вебометрических исследований относятся задачи анализа и выявления гиперссылочных структур различных сегментов Веб-пространства (например, академический сегмент Веба, университетский, и др.), решение которых влияет на качество присутствия этих сегментов в Вебе, на результаты ранжирования поисковых машин (Google, Yandex и др.) или, в случае университетского Веба, на вебометрический рейтинг (http://www.webometrics.info) различных университетов мира \cite{BlekanovSergeevPechnikov}.

 Для получения и обработки больших объемов информации о веб-сайтах и их гиперссылках используются Веб-краулеры (поисковые роботы), общей задачей которых является специализированный обход Веба с целью сбора информации или определения гиперссылочной структуры и полезности каких-либо информационных ресурсов.
 
\textbf{ Эксперимент.} В эксперименте ставилась задача анализа и выявления гиперссылочной структуры Веб-пространства Санкт-Петербургского государственного университета (СПбГУ). 
 
 Для эксперимента использовался программный комплекс обобщенного ядра поискового робота, который обладает высокой гибкостью и масштабируемостью в сравнении с зарубежными аналогами, сильно уступающими в производительности собора и обработки веб-ресурсов и имеющими слабую приспособленность к анализу российского сегмента Веба \cite{BlekanovSergeevMartynenko}. 
 
 К Веб-краулеру с обобщенным ядром дополнительно был разработан и добавлен специализированный алгоритм обхода веб-страниц, который собирает и обрабатывает только страницы из Веб-пространства СПбГУ. В свою очередь пространство СПбГУ состоит из веб-сайта главного домена и сайтов всех его поддоменов (Рис.~\cref{fig:spbuWebSpace}).
 
\begin{figure}[ht]
    \centerfloat{
	        \includegraphics[scale=0.5]{spbuWebSpace}
	    }
    \caption{Веб-пространство СПбГУ.}\label{fig:spbuWebSpace}
\end{figure}

Используя программный комплекс на основе обобщенного ядра поискового робота со специализированным алгоритмом, запущенного с начального множества веб-страниц, требовалось в автоматизированном режиме получить значения следующих показателей, характеризующих гиперссылочную структуру Веб-пространства СПбГУ:
\begin{itemize}
	\item объем Веб-пространства СПбГУ (количество всех различных веб-страниц из Веб-пространства СПбГУ);
	\item количество всех поддоменов из Веб-пространства СПбГУ; 
	\item количество тупиковых (не имеющих ссылок) веб-страниц; 
	\item количество неработающих гиперссылок; 
	\item количество гиперссылок на внешние веб-ресурсы; 
	\item количество поддоменов, связанных с «Главной страницей» главного домена; 
	\item количество поддоменов, несвязанных с «Главной страницей» главного домена; 
	\item количество страниц, имеющие гиперссылки на «Главную страницу» главного домена; 
	\item количество страниц, не имеющие гиперссылки на «Главную страницу» главного домена; 
	\item гиперссылочная структура Веб-пространства СПбГУ в виде матрицы смежности. 
\end{itemize}

В качестве начального множества веб-страниц, с которого Веб-краулер запускал процесс сбора и обработки веб-ресурсов, брался URL-адрес главного веб-сайта СПбГУ -- «http://www.spbu.ru/».

\textbf{Результаты эксперимента.} В ходе эксперимента всего Веб-краулером было обработано и проанализировано \textit{6\space429\space963} гиперссылки, которые содержались на страницах Веб-пространства СПбГУ. Из них: объем ссылок на внешние источники информации равен \textit{507\space168}, а объем внутренних ссылок (на страницы главного домена сайта СПбГУ и его поддоменов) -- \textit{5\space922\space795}. Кроме того, были получены следующие результаты (Табл.~\cref{tab:indicators}):

\begin{table} [htbp]%
	\centering
	\caption{}%
	\label{tab:indicators}% label всегда желательно идти после caption
	\renewcommand{\arraystretch}{1.5}%% Увеличение расстояния между рядами, для улучшения восприятия.
	\begin{SingleSpace}
		\begin{tabulary}{\textwidth}{@{}>{\zz}L >{\zz}C@{}} %Вертикальные полосы не используются принципиально, как и лишние горизонтальные (допускается по ГОСТ 2.105 пункт 4.4.5) % @{} позволяет прижиматься к краям
			\toprule     %%% верхняя линейка
			Показатель & Значение показателя  \\
			\midrule %%% тонкий разделитель. Отделяет названия столбцов. Обязателен по ГОСТ 2.105 пункт 4.4.5
			объем Веб-пространства СПбГУ & 71688 \\				
			количество всех поддоменов из Веб-пространства СПбГУ & 315 \\
			количество поддоменов, сильно-связанных с «Главной страницей» главного домена & 280 \\
			количество поддоменов, несвязанных с «Главной страницей» главного домена & 35 \\
			количество тупиковых веб-страниц & 12516 \\
			количество неработающих гиперссылок & 680 \\
			количество страниц, имеющие гиперссылки на «Главную страницу» главного домена & 10632 \\
			количество страниц, не имеющие гиперссылки на «Главную страницу» главного домена & 61056 \\
			\bottomrule %%% нижняя линейка
		\end{tabulary}%
	\end{SingleSpace}
\end{table}

Также стоит отметить, что для простоты исследований и получения численных значений показателей было предложено хранить выявленную гиперссылочную структуру Веб-пространства СПбГУ в виде матрицы смежности A размерностью [71688\(\times\)71688]. С помощью нее и ее транспонированной формы удобно определять связанность между страницами и в дальнейшем оптимизировать структуру ссылок. 

\textbf{Общие выводы.} По результатам эксперимента видно, что \textit{17,5\%} веб-страниц из всех страниц Веб-пространства СПбГУ занимают тупиковые страницы, которые, например, не учитывает поисковая машина Google в своем алгоритме ранжирования PageRank, что негативно влияет на коммуникабельность всего Веб-пространство университета в целом. 

Также из эксперимента наблюдается плохая связанность страниц Веб-пространства СПбГУ и его поддоменов со страницами главного сайта университета: только \textit{14,8\%} всех веб-страниц имеет ссылки на страницы главного сайта, а \textit{35} доменов со всеми своими страницами и вовсе никак не связаны с ним. 

Стоит отметить и большое количество неработающих гиперссылок (\textit{680} уникальных гиперссылок), которые неоднократно повторяются по всему Веб-пространству СПбГУ, тем самым снижают доверие пользователей к ресурсам сайтов университета. 

Таким образом, проведенный эксперимент демонстрирует слабую связанность и коммуникабельность внутренних ресурсов Веб-пространства СПбГУ, что влечет и его слабую позицию в вебометрическом рейтинге сайтов университетов мира.

\section{Анализ гиперссылочной структуры ресурсов}\label{sec:markup}

%В шаблоне для диссертации и автореферата заданы команды рецензирования.
%Они видны при компиляции шаблона в режиме черновика или при установке
%соответствующей настройки (\verb+showmarkup+) в~файле \verb+common/setup.tex+.
%
%Команда \verb+\todo+ отмечает текст красным цветом.
%\todo{Например, так.}
%
%Команда \verb+\note+ позволяет выбрать цвет текста.
%\note{Чёрный, } \note[red]{красный, } \note[green]{зелёный, }
%\note[blue]{синий.} \note[orange]{Обратите внимание на ширину и расстановку
%    формирующихся пробелов, в~результате приведённой записи (зависит также
%    от~применяемого компилятора).}
%
%Окружение \verb+commentbox+ также позволяет выбрать цвет.
%
%\begin{commentbox}[red]
%    Красный текст.
%
%    Несколько параграфов красного текста.
%\end{commentbox}
%
%\begin{commentbox}[blue]
%    Синяя формула.
%
%    \begin{equation}
%        \alpha + \beta = \gamma
%    \end{equation}
%\end{commentbox}
%
%\verb+commentbox+ позволяет закомментировать участок кода в~режиме чистовика.
%Чтобы убрать кусок кода для всех режимов, можно использовать окружение
%\verb+comment+.
%
%\begin{comment}
%Этот текст всегда скрыт.
%\end{comment}
%
%\section{Анализ текстового контента ресурсов}\label{sec:acronyms}
%
%С помощью пакета \texttt{nomencl} можно создавать удобный сортированный список
%сокращений и условных обозначений во время написания текста. Вызов
%\verb+\nomenclature+ добавляет нужный символ или сокращение с~описанием
%в~список, который затем печатается вызовом \verb+\printnomenclature+
%в~соответствующем разделе.
%Для того, чтобы эти операции прошли, потребуется дополнительный вызов
%\verb+makeindex -s nomencl.ist -o %.nls %.nlo+ в~командной строке, где вместо
%\verb+%+ следует подставить имя главного файла проекта (\verb+dissertation+
%для этого шаблона).
%Затем потребуется один или два дополнительных вызова компилятора проекта.
%\begin{equation}
%    \omega = c k,
%\end{equation}
%где \( \omega \) "--- частота света, \( c \) "--- скорость света, \( k \) "---
%модуль волнового вектора.
%\nomenclature{\(\omega\)}{частота света\nomrefeq}
%\nomenclature{\(c\)}{скорость света\nomrefpage}
%\nomenclature{\(k\)}{модуль волнового вектора\nomrefeqpage}
%Использование
%\begin{verbatim}
%\nomenclature{\(\omega\)}{частота света\nomrefeq}
%\nomenclature{\(c\)}{скорость света\nomrefpage}
%\nomenclature{\(k\)}{модуль волнового вектора\nomrefeqpage}
%\end{verbatim}
%после уравнения добавит в список условных обозначений три записи.
%Ссылки \verb+\nomrefeq+ на последнее уравнение, \verb+\nomrefpage+ "--- на
%страницу, \verb+\nomrefeqpage+ "--- сразу на~последнее уравнение и~на~страницу,
%можно опускать и~не~использовать.
%
%Группировкой и сортировкой пунктов в списке можно управлять с~помощью указания
%дополнительных аргументов к команде \verb+nomenclature+.
%Например, при вызове
%\begin{verbatim}
%\nomenclature[03]{\( \hbar \)}{постоянная Планка}
%\nomenclature[01]{\( G \)}{гравитационная постоянная}
%\end{verbatim}
%\( G \) будет стоять в списке выше, чем \( \hbar \).
%Для корректных вертикальных отступов между строками в описании лучше
%не~использовать многострочные формулы в~списке обозначений.
%
%\nomenclature{%
%    \( \begin{rcases}
%        a_n \\
%        b_n
%    \end{rcases} \)%
%}{коэффициенты разложения Ми в дальнем поле соответствующие электрическим и
%    магнитным мультиполям}
%\nomenclature[a\( e \)]{\( {\boldsymbol{\hat{\mathrm e}}} \)}{единичный вектор}
%\nomenclature{\( E_0 \)}{амплитуда падающего поля}
%\nomenclature{\( j \)}{тип функции Бесселя}
%\nomenclature{\( k \)}{волновой вектор падающей волны}
%\nomenclature{%
%    \( \begin{rcases}
%        a_n \\
%        b_n
%    \end{rcases} \)%
%}{и снова коэффициенты разложения Ми в дальнем поле соответствующие
%    электрическим и магнитным мультиполям. Добавлено много текста, так что
%    описание группы условных обозначений значительно превысило высоту этой
%    группы...}
%\nomenclature{\( L \)}{общее число слоёв}
%\nomenclature{\( l \)}{номер слоя внутри стратифицированной сферы}
%\nomenclature{\( \lambda \)}{длина волны электромагнитного излучения в вакууме}
%\nomenclature{\( n \)}{порядок мультиполя}
%\nomenclature{%
%    \( \begin{rcases}
%        {\mathbf{N}}_{e1n}^{(j)} & {\mathbf{N}}_{o1n}^{(j)} \\
%        {\mathbf{M}_{o1n}^{(j)}} & {\mathbf{M}_{e1n}^{(j)}}
%    \end{rcases} \)%
%}{сферические векторные гармоники}
%\nomenclature{\( \mu \)}{магнитная проницаемость в вакууме}
%\nomenclature{\( r, \theta, \phi \)}{полярные координаты}
%\nomenclature{\( \omega \)}{частота падающей волны}
%
%С помощью \verb+nomenclature+ можно включать в~список сокращения,
%не~используя их~в~тексте.
%% запись сокращения в список происходит командой \nomenclature,
%% а не употреблением самого сокращения
%\nomenclature{FEM}{finite element method, метод конечных элементов}
%\nomenclature{FIT}{finite integration technique, метод конечных интегралов}
%\nomenclature{FMM}{fast multipole method, быстрый метод многополюсника}
%\nomenclature{FVTD}{finite volume time-domain, метод конечных объёмов
%    во~временной области}
%\nomenclature{MLFMA}{multilevel fast multipole algorithm, многоуровневый
%    быстрый алгоритм многополюсника}
%\nomenclature{BEM}{boundary element method, метод граничных элементов}
%\nomenclature{CST MWS}{Computer Simulation Technology Microwave Studio
%    программа для компьютерного моделирования уравнен Максвелла}
%\nomenclature{DDA}{discrete dipole approximation, приближение дискретиных
%    диполей}
%\nomenclature{FDFD}{finite difference frequency domain, метод конечных
%    разностей в~частотной области}
%\nomenclature{FDTD}{finite difference time domain, метод конечных разностей
%    во~временной области}
%\nomenclature{MoM}{method of moments, метод моментов}
%\nomenclature{MSTM}{multiple sphere T-Matrix, метод Т-матриц для множества
%    сфер}
%\nomenclature{PSTD}{pseudospectral time domain method, псевдоспектральный метод
%    во~временной области}
%\nomenclature{TLM}{transmission line matrix method, метод матриц линий передач}
%
\FloatBarrier
